若分析停在"鱼类资源量增加"，2025年恢复至1.8倍这一观测事实已经足够；问题一还要求解释鱼类为何增长，并判断基础资源的承载压力是否改变。将四大家鱼简化为单一消费者功能群、将水草、浮游植物、浮游动物和底栖饵料汇总为单一资源时，总量可以对上，资源层级之间的差别却被抹去，因而难以说明哪一层先承压，也难以联系"水下荒漠"等底层生态退化现象。为保留这部分解释能力，模型按基本食性对应关系分别设置四类基础资源和四大家鱼。
